\section{研究计划进度安排及预期目标}

\subsection{进度安排}
2026年3月25日：提交文献综述，开题报告，文献翻译。

2026年4月1日：完成三个处理器的 SoC 的构建，编写模拟外设设备；使用 riscv-test 测试 SoC 功能正确性。

2026年4月8日：在已有的 SoC 上复现现有的侧信道漏洞。对复现成功/复现失败的结果分析原因。

2026年4月15日：对接生成微架构状态设置程序框架，并测试功能正确性。

2026年4月22日：完成 SoC 时序单元列表导出功能。

2026年4月29日：完成波形差分对齐功能，并与微架构状态设置程序联合测试。

2026年5月6日：完成形式化验证框架从抽象原件，建立假设，设置安全属性的自动化流程，开始进行漏洞检测。

2026年5月13日：分析漏洞检测结果与漏洞检测效率，向开源处理器开发者做漏洞披露报告，进一步完善漏洞检测框架。

2026年5月18号：提交毕业论文。


\subsection{预期目标}

实现 RISC-V 处理器侧信道漏洞检测框架。
该框架能够检测出大部分已有的代表性处理器侧信道漏洞，并且发现新的高质量，可利用的新的侧信道漏洞。
同时相比已有的处理器侧信道漏洞挖掘框架，拥有更高的效率。